\scanpage{30}\mspage{29}
% L1–3
\deleted{\uline{Définition}}

Soit $\operatorname{Tr}(L)$ l'ens \deleted{ordonné} des ps. tronçons de $L$, ordonné par $\ll$.

% L4–6
Pour toute partie finie tot. ordonnée $\Phi$ de $\operatorname{Tr}(L)$, considérons
\[
\deleted{U}X_{\Phi}=\bigcup_{i\in\Phi}T_i\qquad\text{(NB réunion disjointe)}
\]
% L7–12
% slanted
\margin{\uline{Prop}}
Alors l'application
\diagram{\begin{tikzcd}[ampersand replacement=\&,column sep=large,row sep=1.3ex]
\Phi \arrow[r,mapsto] \& \deleted{U}X_{\Phi}\\
\underset{\substack{\text{drapeaux de}\\\operatorname{Tr}(L),\ \text{y inclus}\\\text{\unsure{par|pour|pav|pau}}}}
{\underline{{\operatorname{Drap}}_{\ll}^{\unsure{\partial|b|*|\varnothing}}(\operatorname{Tr}(L))}}
\arrow[r] \& \unsure{\mathcal P|\mathscr P|\mathrm P}(L)
\end{tikzcd}}
est \uline{injective}.

% L13–15
Les ens de la forme $U_{\Phi}$ s'appellent \uline{modérés}. Pour un tel
\[
\deleted{U}X_{\Phi}=\bigcup T_i
\]
on pose
% L16–17
\[
\operatorname{Fr}(\deleted{U}X_{\Phi})=
\underbrace{\left(\bigcup_i\partial T_i\right)}_{\substack{\deleted{\text{\unsure{est}}}\\\partial U,\ \text{drapeau de } L}}\setminus\varepsilon
\]
% slanted
\margin{``frontière'' relative dans $L$ (mais pas pour la topologie \unsure{de|sur} \unsure{m.|sur} $L$\ldots)}
% L18–24
où
\[
\begin{aligned}
\varepsilon={}&\bigl\{t\in\partial L\cap U\mid
  \deleted{\unsure{t\in U}}\ \{t\}\notin\Phi\quad\text{i.e. }\deleted{\unsure{\Phi|\notin}}\\
&\hspace{6em}\{t\}\text{ n'est pas parmi les }T_i,\\
&\hspace{6em}\deleted{\text{i.e. }\exists\,x\in L,\ \text{\unsure{tq} } x<{}}\\
&\hspace{6em}\deleted{\text{i.e. il y}}\quad\text{i.e }U\text{ est un}\\
&\hspace{8em}\text{voisinage de }t\text{ dans } L\\
&\hspace{6em}\text{i.e. }\{t\}\text{ pas point isolé}\\
&\hspace{9em}\text{de }U\ \ldots\bigr\}
\end{aligned}
\]
% slanted
% L lower left note
\margin{Les $T_i$ s'appellent les ``composantes (connexes!) de \deleted{\unsure{$U$}}$X$ (mais pas pour la top canonique) \unsure{l'on|l'ens} \unsure{va|en} \unsure{aussi|sera} \unsure{les|se} \unsure{noter|noté} parfois $\Pi_{\unsure{0|o}}(X)$}
